%==============================================================================
%  CAPÍTULO I
%==============================================================================
\chapter{Evolución histórica y dogmática del derecho concursal chileno}
\label{cap:evolucion}

\fichaCapitulo{De la punición a la reorganización.}{Transición sistémica del sistema
concursal, consagración de los principios rectores contemporáneos y crisis de la noción
formal de insolvencia.}

%==============================================================================
\bloqueTeorico
%==============================================================================

\subsection{El desplazamiento del eje: de la sanción al salvataje}

El derecho concursal chileno experimentó, en el curso de un siglo y medio, un
desplazamiento completo de su centro de gravedad. El sistema que instauró el \ccom{} de
1865\nocite{codigo-comercio} y que perfeccionaron sucesivamente la \ley{4.558}
\citeyearpar{ley-4558} y la \ley{18.175} \citeyearpar{ley-18175} concebía la quiebra
como un episodio esencialmente patológico y reprochable: el comerciante que
cesaba en el pago de sus obligaciones era, ante todo, un sujeto sospechoso, y el
procedimiento que se abría en su contra tenía por función simultánea liquidar su
patrimonio y depurar el mercado de su presencia. La \LIR{}, en cambio, parte de una
premisa inversa. La insolvencia es un riesgo ordinario e inevitable de la actividad
económica \parencite{aj-texto-refundido}, y la función del ordenamiento no consiste en
castigar al que la sufre, sino en administrar sus consecuencias de manera que se preserve el valor todavía existente y
se permita el retorno del deudor honesto al circuito productivo.

Este capítulo inicial sienta las bases conceptuales indispensables para comprender esa
transición. No se trata de una mera actualización terminológica ---la sustitución de la
voz «quiebra» por la de «liquidación», o la del «síndico» por el «liquidador»---, sino
de un cambio en la función que el ordenamiento asigna al procedimiento colectivo. Bajo
el régimen anterior, la quiebra operaba como una sanción civil, y en ciertos supuestos
cuasipenal, impuesta al comerciante insolvente. Sus efectos --- la inhabilitación, el
estigma registral, la presunción de culpabilidad en la calificación --- estaban
diseñados para disuadir la insolvencia antes que para gestionarla.

La consecuencia práctica de ese diseño fue paradójica y merece un análisis detenido, pues
explica buena parte de las decisiones de política legislativa que informan tanto la
\LIR{} como su reforma. Al elevar el costo reputacional y patrimonial de acogerse al
procedimiento colectivo, el sistema clásico desincentivaba precisamente la conducta que
habría permitido conservar valor: el reconocimiento temprano de la crisis. El deudor
racional postergaba el ingreso al concurso hasta que la empresa había consumido no sólo
su patrimonio, sino también su capital reputacional, su cartera de clientes y su equipo
humano. Cuando finalmente se abría la quiebra, el activo realizable era una fracción
marginal del que existía cuando la crisis se manifestó. Los acreedores recuperaban
menos; los trabajadores perdían empleos que eran viables; y la economía en su conjunto
soportaba una destrucción de valor que no era consecuencia de la insolvencia, sino del
retardo en su tratamiento.

\begin{practica}[Por qué esto importa en la asesoría]
La lectura histórica no es ornamental. El profesional que asesora a un deudor en crisis
enfrenta, todavía hoy, la resistencia cultural que aquel sistema sembró: el cliente
percibe el concurso como una confesión de fracaso. Explicar que la \LIR{} invirtió esa
lógica ---y que el ingreso temprano al procedimiento amplía, no reduce, las opciones
disponibles--- forma parte del núcleo de la asesoría concursal contemporánea.
\end{practica}

\subsection{Justificación de la extensión del capítulo}

La extensión asignada a este capítulo se justifica en la necesidad de analizar de forma
exhaustiva dos cuerpos de materiales que la práctica suele dar por conocidos sin serlo.

El primero es la doctrina clásica de la quiebra. Su estudio no responde a un interés
anticuario: numerosas categorías de la \LIR{} ---la cesación de pagos, el
\desasimiento{}, la masa, el período sospechoso--- fueron construidas por la dogmática
del siglo \textsc{xx} y conservan hoy el contenido que aquélla les dio. Quien las aplica
sin conocer su genealogía corre el riesgo de atribuirles un alcance que el legislador de
2014 deliberadamente modificó, o de desconocer el que deliberadamente mantuvo.

El segundo es el modo en que los principios de la \LIR{} redefinieron las relaciones de
crédito en el país. La reforma concursal no operó en el vacío: alteró la posición
relativa de los acreedores financieros, de los proveedores, de los trabajadores y del
fisco, y lo hizo mediante reglas cuya justificación sólo resulta inteligible desde los
principios que las informan.

%==============================================================================
\bloqueNormativo
%==============================================================================

El capítulo se organiza en torno a cuatro bloques temáticos, que se desarrollan en las
secciones siguientes.

\subsection{La evolución legislativa chilena}

\subsubsection{El Código de Comercio de 1865 y el régimen de la quiebra mercantil}

El \ccom{} de 1865 reguló la quiebra en su Libro \textsc{iv}, siguiendo de cerca el
modelo francés de 1807 y su tradición de raíz estatutaria italiana. Tres rasgos
caracterizan ese régimen originario y conviene fijarlos con precisión, porque cada uno
de ellos fue objeto de reforma posterior.

\begin{enumerate}
  \item \textbf{Carácter mercantil del presupuesto subjetivo.} La quiebra era un
  instituto reservado al comerciante. El deudor civil quedaba sujeto al régimen común de
  la ejecución individual y a la cesión de bienes del \cc{}, sin acceso a un
  procedimiento colectivo estructurado.

  \item \textbf{Carácter sancionatorio del procedimiento.} La declaratoria producía la
  inhabilitación del fallido y abría la calificación de la quiebra, orientada a
  determinar si ésta era fortuita, culpable o fraudulenta, con las consecuencias penales
  correspondientes.

  \item \textbf{Vocación exclusivamente liquidatoria.} El convenio existía, pero como
  excepción negociada al curso natural del procedimiento, que era la realización de los
  bienes y el reparto del producto.
\end{enumerate}

\subsubsection{La Ley \Nro 4.558 y la institucionalización de la sindicatura}

La \ley{4.558} \citeyearpar{ley-4558} introdujo un cambio de arquitectura institucional
cuya importancia
excede su contenido normativo: profesionalizó la administración del concurso mediante la
creación de la Sindicatura General de Quiebras y sometió al síndico a un régimen de
control público.\verificar{Fecha exacta de publicación y ámbito de la Ley 4.558.} Con
ello se reconoció, por primera vez de manera explícita, que la administración de un
patrimonio en crisis es una función de interés general y no un asunto privado entre
deudor y acreedores.

\subsubsection{La Ley \Nro 18.175 y su incorporación al Código de Comercio}

La \ley{18.175} \citeyearpar{ley-18175} constituyó el estatuto concursal aplicable en
Chile durante más de tres décadas. Su rasgo distintivo fue el desplazamiento parcial hacia un modelo de
administración privada de la quiebra, con la sustitución de la sindicatura estatal por un
sistema de síndicos privados sometidos a fiscalización.\verificar{Precisar la evolución
de la Superintendencia de Quiebras y su relación con la Fiscalía Nacional de Quiebras.}
Posteriormente, su articulado fue incorporado como Libro \textsc{iv} del \ccom{},
consolidando formalmente la naturaleza mercantil del instituto.\verificar{Confirmar la
ley que incorporó el texto de la Ley 18.175 como Libro IV del Código de Comercio y su
año.}

\subsubsection{La irrupción de la Ley \Nro 20.720}

La \LIR{} \citeyearpar{ley-20720} derogó el régimen anterior y sustituyó la lógica del
sistema. Siguiendo la sistematización de \textcite{aj-texto-refundido}, sus innovaciones
estructurales pueden agruparse en cinco:

\begin{enumerate}
  \item \textbf{Superación del presupuesto mercantil.} La ley abandona la distinción
  entre deudor comerciante y no comerciante, y la reemplaza por la dicotomía entre
  \empresadeudora{} y \personadeudora{}, construida sobre criterios tributarios y no
  sobre la naturaleza de los actos ejecutados.

  \item \textbf{Paridad estructural entre reorganización y liquidación.} La
  reorganización deja de ser un incidente del procedimiento liquidatorio y pasa a
  constituir un procedimiento autónomo, con presupuestos, plazos y efectos propios.

  \item \textbf{Institucionalización de la \Superir{}.} Se crea un órgano fiscalizador
  con competencias normativas, de registro y sancionatorias sobre veedores y
  liquidadores.

  \item \textbf{Publicidad concentrada en el \boletin{}.} Se sustituye la publicación
  fragmentaria en el Diario Oficial y en diarios de circulación por una plataforma
  electrónica única, gratuita y de acceso público.

  \item \textbf{Consagración del \discharge{}.} Se introduce la extinción de los saldos
  insolutos del deudor persona natural, rompiendo con el principio de responsabilidad
  patrimonial universal e ilimitada del \artcc{2465} \citeyearpar{codigo-civil}. Sobre el
  alcance de esta ruptura, véase \textcite{doc-discharge-modelo}.
\end{enumerate}

\subsubsection{La reforma de la Ley \Nro 21.563}

La \Lreforma{} \citeyearpar{ley-21563} no sustituyó el sistema de 2014, sino que corrigió
sus disfunciones operativas \parencite{aj-ley21563, aj-reporte-21563}. El diagnóstico
legislativo que la precedió identificó un problema central:
los procedimientos diseñados para la empresa mediana y grande resultaban
económicamente inaccesibles para las micro y pequeñas empresas y para las personas
naturales, de modo que el sistema atendía sólo a una fracción menor de la insolvencia
efectivamente existente. La respuesta consistió en la creación de procedimientos
simplificados, la reducción de exigencias documentales y la introducción de un control
de probidad ---el incidente de mala fe--- que compensara la menor intensidad de
verificación previa \parencite{aj-prado-modificaciones}. Sobre las brechas que la reforma
no alcanzó a cerrar respecto de las empresas de menor tamaño, véase
\textcite{prensa-df-brechas-mypes}. Estas materias se desarrollan en los capítulos
\ref{cap:estatuto-deudores}, \ref{cap:reorganizacion-simplificada},
\ref{cap:liquidacion-simplificada} y \ref{cap:mala-fe}.

\subsection{Los principios rectores contemporáneos}

\subsubsection{Conservación de la empresa viable}

El principio de conservación no ordena mantener con vida a toda empresa, sino distinguir
entre la empresa económicamente viable que atraviesa una crisis financiera y la empresa
estructuralmente inviable. Respecto de la primera, la liquidación destruye valor: el
conjunto organizado de bienes vale más operando que desmembrado. Respecto de la segunda,
la conservación lo destruye igualmente, al consumir recursos en sostener una actividad
que no genera valor.

La consecuencia dogmática es relevante: el principio de conservación no es un mandato de
resultado dirigido al juez, sino un criterio de asignación procedimental. Impone que el
sistema ofrezca una vía de reorganización accesible y que la liquidación opere como
\lat{ultima ratio}, pero no autoriza a imponer la continuidad de una empresa cuya
inviabilidad los acreedores han constatado.

\subsubsection{Protección de la masa y \lat{par condicio creditorum}}

La \lat{par condicio creditorum} ---el trato paritario de los acreedores de igual
rango--- constituye el fundamento último del carácter colectivo del procedimiento. Si
cada acreedor pudiera perseguir individualmente los bienes del deudor insolvente, el
resultado sería una carrera en que el más diligente o el mejor informado se satisface
íntegramente mientras los restantes nada obtienen, sin que esa distribución responda a
criterio alguno de justicia o de eficiencia.

Conviene precisar el alcance del principio: no impone igualdad absoluta, sino igualdad
dentro de cada categoría. El sistema de preferencias del \artcc{2472} \citeyearpar{codigo-civil} subsiste
íntegramente en el concurso, y el propio legislador ha reforzado algunas de ellas ---
señaladamente las de origen laboral, según se examina en el capítulo
\ref{cap:laboral-penal} \parencite{doc-proteccion-trabajadores}.

\subsubsection{Universalidad subjetiva y objetiva}

La universalidad opera en dos planos. En el plano \emph{objetivo}, el procedimiento
comprende la totalidad de los bienes del deudor, con la sola excepción de los legalmente
inembargables. En el plano \emph{subjetivo}, comprende a la totalidad de los acreedores,
que quedan sujetos a sus resultados hayan o no verificado sus créditos.

La tensión más aguda que este principio enfrenta en el derecho vigente se produce a
propósito de las deudas que pretenden sustraerse a sus efectos por vía de especialidad
normativa. El caso paradigmático ---el Crédito con Aval del Estado--- se estudia en el
capítulo \ref{cap:discharge-cae}.

\subsubsection{Celeridad procesal}

La celeridad no responde a una preferencia estética por los procedimientos breves, sino a
una constatación económica: el valor del activo concursal decrece con el tiempo. Las
manifestaciones concretas del principio en la \LIR{} son numerosas ---plazos fatales,
notificación por el \boletin{}, régimen recursivo restrictivo, audiencias concentradas---
y varias de ellas han sido objeto de reproche constitucional, según se examina en el
capítulo \ref{cap:liquidacion-ordinaria}.

\subsection{La noción de insolvencia y sus categorías vecinas}

El ordenamiento chileno no define la insolvencia. La \LIR{} \citeyearpar{ley-20720} la
presupone, la emplea como criterio de apertura y la asocia a hipótesis de hecho tasadas, pero no la conceptualiza.
Esta omisión ---compartida con otros ordenamientos--- ha obligado a la doctrina y a la
jurisprudencia a construir la noción, y explica la persistencia del debate que se examina
en la sección \ref{sec:cap01-doctrina}.

Para el trabajo forense resulta indispensable distinguir cuatro estados que el lenguaje
corriente confunde:

\begin{description}[leftmargin=0pt,style=nextline,font=\sffamily\bfseries]
  \item[Iliquidez] Imposibilidad transitoria de atender obligaciones exigibles por falta
  de disponibilidad inmediata de medios de pago, subsistiendo un activo cuyo valor
  excede al pasivo. Es un problema de calce temporal entre flujos, no de suficiencia
  patrimonial.

  \item[Cesación de pagos] Manifestación externa del incumplimiento generalizado. Es un
  hecho revelador, no un estado patrimonial: puede existir cesación sin insolvencia
  (deudor solvente que se niega a pagar) e insolvencia sin cesación (deudor deficitario
  que se financia con crédito nuevo).

  \item[Sobreendeudamiento] Desproporción entre el nivel de endeudamiento y la capacidad
  de generación de ingresos del deudor, típicamente referida a la persona natural. El
  patrimonio puede ser positivo, pero el servicio de la deuda absorbe una porción
  insostenible del ingreso.

  \item[Insolvencia] Estado patrimonial deficitario y estructural que impide al deudor
  cumplir de forma íntegra y oportuna sus obligaciones exigibles.
\end{description}

\begin{alerta}[Consecuencia práctica de la distinción]
La calificación no es teórica: determina la vía procedimental idónea. La iliquidez
admite tratamiento reorganizatorio; la insolvencia estructural rara vez lo tolera, y la
propuesta de acuerdo que la desconoce está destinada al rechazo de los acreedores o,
peor, a un cumplimiento incumplido que consumirá el activo remanente.
\end{alerta}

\subsection{La justificación económica del \discharge{}}

La extinción de los saldos insolutos constituye la innovación de mayor densidad
valorativa de la \LIR{}. Su justificación admite tres lecturas complementarias.

\paragraph{Lectura de eficiencia.} La deuda perpetua e impagable no genera recuperación
para el acreedor y sí destruye incentivos en el deudor: quien sabe que todo ingreso
adicional será absorbido por acreedores preexistentes carece de motivo racional para
generarlo, y se desplaza hacia la informalidad. El \discharge{} restituye el incentivo
marginal al esfuerzo productivo.

\paragraph{Lectura de asignación de riesgo.} El acreedor profesional está en mejor
posición que el deudor para evaluar, diversificar y tarificar el riesgo de crédito. La
descarga asigna el riesgo residual a quien puede administrarlo a menor costo.

\paragraph{Lectura de dignidad.} La responsabilidad patrimonial universal del
\artcc{2465}, aplicada sin límite temporal a la persona natural, convierte una deuda en
una condición vital. El \freshstart{} expresa la decisión de que el fracaso económico no
determine irreversiblemente la biografía de una persona.

\medskip
Frente a estas razones se alza el argumento del riesgo moral: si la deuda puede
extinguirse, el deudor tiene incentivos para endeudarse en exceso o para provocar su
propia insolvencia. La respuesta del sistema chileno a esa objeción no consiste en
restringir el \discharge{}, sino en condicionarlo a la buena fe del deudor mediante el
incidente del \art{169 A}, introducido por la \Lreforma{} y examinado en el capítulo
\ref{cap:mala-fe}.

\subsection{El sistema chileno en el mapa comparado}

La \LIR{} \citeyearpar{ley-20720} no nació aislada: se inscribe en una corriente mundial de
modernización del derecho de la insolvencia que desplazó el eje desde la liquidación
sancionatoria hacia la reorganización y la segunda oportunidad. Situar el modelo chileno en ese
mapa permite comprender qué tomó de cada tradición y dónde se apartó.

\subsubsection{La tradición estadounidense: reorganización y \emph{fresh start}}

El \emph{Bankruptcy Code} \parencite{us-bankruptcy-code} estadounidense es la matriz de la que
proviene buena parte del vocabulario contemporáneo. Su \destacar{Chapter 11} consagró el modelo
de reorganización con el deudor en posesión (\emph{debtor in possession}), y su \destacar{Chapter
7} la liquidación con \emph{fresh start}. La \LIR{} recogió esa dualidad ---reorganización y
liquidación como procedimientos autónomos--- y el propio concepto de \discharge{}, aunque con un
diseño institucional distinto: en Chile la administración pasa al liquidador, no permanece en el
deudor.

\subsubsection{La tradición europea: la Directiva 2019/1023 y la reestructuración preventiva}

En Europa, la \destacar{Directiva (UE) 2019/1023} \parencite{ue-directiva-2019-1023} impuso a los
Estados miembros marcos de reestructuración \emph{preventiva} ---antes de la insolvencia
declarada--- y de exoneración de deudas para el empresario de buena fe. España la transpuso
mediante la \destacar{Ley 16/2022} \parencite{esp-ley16-2022}, que sustituyó su régimen anterior
por la exoneración del pasivo insatisfecho.

\subsubsection{Qué distingue al modelo chileno}

Del contraste emergen tres rasgos propios del sistema chileno, que esta obra examina en detalle:

\begin{enumerate}
  \item \textbf{Administración por un tercero.} A diferencia del \emph{debtor in possession}
  norteamericano, la liquidación chilena produce \desasimiento{}: administra el liquidador
  (capítulo \ref{cap:liquidacion-ordinaria}).
  \item \textbf{Un procedimiento administrativo para la persona.} La renegociación ante la
  \superir{} ---gratuita y no judicial--- carece de equivalente directo en los modelos de
  referencia (capítulo \ref{cap:renegociacion}).
  \item \textbf{La ausencia de un catálogo de deudas excluidas.} Donde EE.UU. y España fijan en
  la ley qué deudas sobreviven al \discharge{}, Chile lo dejó a los tribunales, con la
  inestabilidad que documenta el capítulo \ref{cap:discharge-cae}.
\end{enumerate}

\begin{foco}[Por qué el comparado importa para litigar en Chile]
El derecho comparado no es ornamento erudito. Cuando la ley chilena calla ---como en la suerte
del crédito educativo o en el estándar de la mala fe--- los tribunales miran cómo resolvieron
los sistemas que Chile tomó como modelo. Conocer el \emph{undue hardship} estadounidense o la
exoneración española da argumentos concretos, no sólo contexto.
\end{foco}

%==============================================================================
\bloquePractico
%==============================================================================

\subsection{El diagnóstico preliminar del cliente en crisis patrimonial}

La primera decisión de la asesoría concursal ---y la que condiciona todas las
posteriores--- consiste en calificar la naturaleza de la crisis. El siguiente flujograma
sintetiza el árbol de decisión que debe recorrerse antes de proponer vía alguna.

\begin{flujograma}{Evaluación inicial del cliente deudor}{cap01-diagnostico}
  \node[hito] (eval) {EVALUACIÓN INICIAL DEL CLIENTE DEUDOR\\[2pt]
    \normalfont Contraste activo realizable / pasivo exigible};

  \coordinate (split) at ($(eval.south)+(0,-6mm)$);
  \node[nodo, below left=12mm and 6mm of eval.south] (temporal)
    {INSOLVENCIA TEMPORAL O ILIQUIDEZ\\[2pt]
     \normalfont Flujos de caja suspendidos, pero activos realizables viables};
  \node[nodo, below right=12mm and 6mm of eval.south] (estructural)
    {INSOLVENCIA ESTRUCTURAL O CRÓNICA\\[2pt]
     \normalfont El pasivo exigible supera con creces el activo liquidable};

  \node[favorable, below=10mm of temporal] (reorg)
    {VÍA DE LA REORGANIZACIÓN\\[2pt]
     \normalfont Reestructuración de pasivos y activos};
  \node[adverso, below=10mm of estructural] (liq)
    {VÍA DE LA LIQUIDACIÓN\\[2pt]
     \normalfont Realización colectiva de activos};

  \draw[flecha] (eval.south) -- (split) -| (temporal.north);
  \draw[flecha] (eval.south) -- (split) -| (estructural.north);
  \draw[flecha] (temporal) -- (reorg);
  \draw[flecha] (estructural) -- (liq);
\end{flujograma}

\subsection{Cuestionario técnico de diagnóstico}

La aplicación del árbol anterior exige información que el cliente rara vez aporta de
forma espontánea y completa. El siguiente cuestionario ordena la indagación:

\begin{enumerate}
  \item \textbf{Presupuesto subjetivo.} ¿Ha sido el deudor contribuyente de primera
  categoría en los veinticuatro meses anteriores? De la respuesta depende su calificación
  como \empresadeudora{} o \personadeudora{} y, con ella, el catálogo íntegro de
  procedimientos disponibles (capítulo \ref{cap:estatuto-deudores}).

  \item \textbf{Magnitud.} ¿Cuál fue el nivel de ventas anuales del último ejercicio?
  El umbral de la \Lpyme{} \citeyearpar{ley-20416} determina el acceso a los
  procedimientos simplificados.

  \item \textbf{Composición del pasivo.} ¿Qué proporción corresponde a acreedores
  financieros, proveedores, obligaciones laborales, deudas previsionales y créditos
  fiscales? La distribución anticipa los quórums alcanzables en una eventual junta.

  \item \textbf{Garantías reales.} ¿Existen hipotecas o prendas vigentes? ¿Sobre bienes
  esenciales para el giro? La respuesta condiciona críticamente la viabilidad de una
  reorganización.

  \item \textbf{Personas relacionadas.} ¿Hay acreedores relacionados con el deudor? Su
  exclusión del cómputo de determinados quórums altera el mapa de mayorías.

  \item \textbf{Actos del bienio anterior.} ¿Se ejecutaron enajenaciones, daciones en
  pago, constituciones de garantía sobre deudas preexistentes o actos gratuitos? Estos
  antecedentes anticipan la exposición a acciones revocatorias (capítulo
  \ref{cap:revocatorias}) y, eventualmente, al incidente de mala fe.

  \item \textbf{Ejecuciones en curso.} ¿Existen juicios ejecutivos notificados? En sede
  de renegociación administrativa, este dato incide directamente en la admisibilidad
  (capítulo \ref{cap:renegociacion}).

  \item \textbf{Deudas de tratamiento especial.} ¿Hay obligaciones derivadas de crédito
  educativo, pensiones alimenticias u otras cuya suerte frente al \discharge{} es
  controvertida? (capítulo \ref{cap:discharge-cae}).
\end{enumerate}

\subsection{Verificación documental obligatoria}

El diagnóstico basado exclusivamente en la declaración del cliente es profesionalmente
inaceptable: la experiencia forense demuestra que el deudor en crisis subestima
sistemáticamente su pasivo y omite obligaciones cuya existencia ha desplazado de su
atención. Antes de diseñar estrategia alguna deben obtenerse y contrastarse, como
mínimo, los siguientes antecedentes:

\begin{description}[leftmargin=0pt,style=nextline,font=\sffamily\bfseries]
  \item[Informe consolidado de deudas] Emitido conforme a la normativa de la \cmf{},
  permite establecer el endeudamiento del sistema financiero formal.
  \verificar{Denominación oficial vigente del informe y vía de obtención.}

  \item[Boletines comerciales] Registran protestos e incumplimientos informados, incluida
  deuda no bancaria que no aparece en el informe anterior.

  \item[Carpeta tributaria electrónica del \sii{}] Acredita el inicio de actividades, la
  categoría tributaria, el nivel de ventas y la existencia de deuda fiscal. Es el
  documento que determina, en definitiva, la calificación subjetiva del deudor.

  \item[Certificados de hipotecas, gravámenes y prohibiciones] Sobre cada inmueble del
  deudor, así como el certificado de anotaciones vigentes de vehículos motorizados.

  \item[Cartola de cotizaciones y saldos previsionales] Incluidos los instrumentos de
  ahorro previsional voluntario, cuya omisión en el inventario ha sido tratada por la
  jurisprudencia como indicio de mala fe (capítulo \ref{cap:liquidacion-simplificada}).
\end{description}

\begin{practica}[Asimetría de información]
La finalidad de esta verificación no es sólo completar el inventario. Es detectar
asimetrías de información \emph{antes} de que las detecte la contraparte: un pasivo
omitido que emerge durante la tramitación destruye la credibilidad de la propuesta de
acuerdo, y un activo omitido abre la puerta al incidente del \art{169 A}.
\end{practica}

%==============================================================================
\bloqueJurisprudencial
\label{sec:cap01-doctrina}
%==============================================================================

\subsection{La ausencia de definición legal y sus consecuencias}

La \LIR{} emplea la voz «insolvencia» en su denominación y en su articulado sin
definirla. La consecuencia es que el contenido del concepto debe extraerse de las
hipótesis de apertura y de la elaboración doctrinaria, con el riesgo de que la aplicación
judicial oscile entre criterios incompatibles.

\subsection{La controversia entre la posición formalista y la sustancialista}

\paragraph{Posición formalista.} Asimila la insolvencia al hecho externo del
incumplimiento. Bastaría acreditar la cesación en el pago de una obligación para tener
por configurado el presupuesto del concurso. Su virtud es la certeza probatoria: el hecho
es objetivo, verificable y no exige incursionar en el balance del deudor. Su defecto es
la sobreinclusión: admite en el procedimiento colectivo a deudores solventes que
simplemente no han pagado, y convierte al concurso en un instrumento de presión
disponible para el acreedor individual.

\paragraph{Posición sustancialista.} Define la insolvencia como un estado patrimonial
deficitario y complejo que impide al deudor cumplir íntegra y oportunamente sus
obligaciones exigibles. El incumplimiento es indicio, no definición. Su virtud es la
correspondencia con la finalidad del sistema; su dificultad, la carga probatoria que
impone y la necesidad de un juicio de valor sobre la situación patrimonial.

\medskip
La \LIR{} adopta una solución de compromiso: construye las causales de liquidación
forzosa sobre hechos externos tasados ---técnica formalista--- pero concede al deudor la
posibilidad de enervarlas acreditando su solvencia, lo que reintroduce el criterio
sustancial por vía de defensa. El análisis detallado de este mecanismo corresponde al
capítulo \ref{cap:liquidacion-ordinaria}.

\subsection{Jurisprudencia esencial}

\paragraph{\sentencia{Corte Suprema}{15.988-2013}{2013} \parencite{juris-cs-15988-2013}.}
Conceptualiza la quiebra y la insolvencia como fenómenos de naturaleza colectiva, que
exigen una mirada de protección general de los acreedores por sobre la lógica de la
ejecución individual. La sentencia conserva vigencia interpretativa bajo la \LIR{}, en
cuanto fija el fundamento del carácter universal del procedimiento.
\verificar{Contrastar los considerandos citados con el texto oficial del fallo.}

\begin{alerta}[Nota metodológica sobre las citas jurisprudenciales]
Las sentencias que se citan a lo largo de esta obra han sido seleccionadas por su valor
doctrinario. Antes de invocarlas en un escrito, el lector debe verificar su vigencia,
la eventual existencia de fallos posteriores en sentido contrario y, tratándose de
materias afectadas por la \Lreforma{}, si el criterio recayó sobre el texto anterior o
posterior a la reforma.
\end{alerta}

\subsection{Bibliografía específica del capítulo}

Para la profundización de las materias tratadas se recomienda, sobre la evolución del
sistema y el texto refundido vigente, \textcite{aj-texto-refundido} y
\textcite{aj-ley20720-evolucion}; sobre el alcance de la reforma,
\textcite{aj-reporte-21563} y \textcite{doc-nuevos-aires}; y sobre el fundamento y los
límites del \discharge{}, \textcite{doc-discharge-modelo}.
